\documentclass[11pt]{article}
\usepackage{amsmath, amssymb, amsthm}
\usepackage[margin=1in]{geometry}
\usepackage{setspace}
\usepackage{booktabs}
\setstretch{1.15}

\newtheorem{result}{Result}[section]
\newtheorem{assumption}{Assumption}

\title{Trade Tariffs and Exchange Rates in a Multilateral World\\ \large Model derivations, restructured}
\author{}
\date{\today}

\begin{document}
\maketitle

\section{Introduction: What Survives from the Two-Country Model?}

A familiar result in international economics is that an import tariff appreciates the currency of the country imposing it. The intuition is clearest in a two-country world. At the initial exchange rate, a tariff reduces import demand and creates a trade surplus for the tariffing country. Provided the Marshall--Lerner condition holds, an appreciation is required to restore trade balance.

In a multilateral world this statement is incomplete. Once there is more than one foreign supplier, a tariff does not only change the tariffing country's demand for foreign goods in aggregate. It also changes the allocation of that demand across foreign countries. Expenditure displaced from the tariffed country can be diverted toward untreated trading partners, changing their trade balances and exchange rates as well.

This creates a distinction that does not exist in the two-country model: the tariffing country may appreciate against the country it tariffs while depreciating against some third countries. The bilateral statement that ``a tariff appreciates the currency'' therefore need not survive multilateralization.

The central result of these notes is that a closely related aggregate statement does. In the symmetric $N$-country model, a unilateral tariff always appreciates the tariffing country's nominal effective exchange rate, even when some individual bilateral exchange rates move in the opposite direction. Cross-origin substitution determines how the exchange-rate adjustment is distributed across foreign partners, but it does not overturn the sign of the aggregate response.

The results can be understood through a distinction that recurs throughout the analysis:
\[
\boxed{\text{aggregate home-versus-foreign adjustment}}
\qquad\text{versus}\qquad
\boxed{\text{reallocation across foreign suppliers}} .
\]
The first margin governs the aggregate value of the tariffing country's currency. The second determines how the adjustment is distributed across individual foreign partners. The two-country model contains only the first observable margin, because there is only one foreign supplier. A multilateral model contains both. Bilateral exchange rates reflect both margins, whereas aggregation across trading partners largely removes the second. In this sense, the conventional two-country result is best understood as an aggregate result that happens also to be bilateral when there is only one foreign country.

The analysis develops this argument progressively. We first recover the conventional two-country result and connect it exactly to the Marshall--Lerner condition (Section 3). We then add a single untreated country (Section 4). This is the smallest environment in which foreign-supplier reallocation exists, and it is enough to generate the central bilateral ambiguity: the tariffing country always appreciates against the target but can depreciate against the untreated country when cross-origin substitution is sufficiently strong. Only after these mechanisms are visible do we introduce the general disturbance-adjustment framework (Section 5), which separates the trade-balance disturbance created by the tariff from the exchange-rate adjustment that eliminates it. Section 6 asks what happens to the overall value of the tariffing country's currency and shows that the bilateral ambiguity disappears from the nominal effective exchange rate. Section 7 generalizes to $N$ countries: increasing the number of countries dilutes trade diversion toward any individual bystander, but leaves the sign of the effective response unchanged under symmetry. Sections 8--10 examine tariff configurations, real exchange rates, and asymmetry. Section 11 relates the mechanism to the canonical theory of preferential trade agreements, in which the familiar Viner--Mundell result emerges as a special environment where the forces opposing trade diversion toward the untreated exporter have been removed. Section 12 concludes. Every closed form is verified numerically in the replication package.

\section{Economic Environment: Two Margins of Substitution}

\subsection{Production and prices}

Consider $N$ countries indexed by $i, j \in \{1, \dots, N\}$. Throughout, $A$ is the tariff imposer, $B$ the target, and $C$ a bystander. Each country produces a country-specific tradable $T_i$ (only $i$ produces it, all countries consume it) and a non-tradable $N_i$, from labor alone:
\begin{equation}
Y_{T_i} = A_{T_i} L_{T_i}, \qquad Y_{N_i} = A_{N_i} L_{N_i}, \qquad L_{T_i} + L_{N_i} = L_i .
\end{equation}
Perfect competition and the normalization $P_{T_i} = 1$ (producer price in $i$'s own currency) give
\begin{equation}
w_i = A_{T_i}, \qquad P_{N_i} = \frac{w_i}{A_{N_i}} = \frac{A_{T_i}}{A_{N_i}} .
\label{eq:wage}
\end{equation}
The production side is deliberately simple. With linear technology and perfect competition, producer prices are pinned down by unit labor costs, and wages are constant in own currency. All relative adjustment therefore runs through exchange rates. Because wages are fixed in domestic currency, movements in nominal exchange rates, common-currency relative wages, and inverse terms of trade coincide one-for-one; the formal statement is given with the definitions in Section 2.4.

\subsection{Preferences: home versus foreign and foreign versus foreign}

The demand side contains the distinction that drives the results. Consumption has two nested margins:
\begin{align}
C_i &= C_{T_i}^{\alpha_T}\, C_{N_i}^{1-\alpha_T}
&& \text{(Cobb--Douglas outer, tradable share } \alpha_T)
\label{eq:outer}\\[2pt]
C_{T_i} &= \Big[\, \alpha_D^{1/\eta}\, C_{T_i i}^{\frac{\eta-1}{\eta}}
+ (1-\alpha_D)^{1/\eta}\, M_i^{\frac{\eta-1}{\eta}} \Big]^{\frac{\eta}{\eta-1}}
&& \text{(home vs.\ imports, elasticity } \eta, \text{ home weight } \alpha_D)
\label{eq:mid}\\[2pt]
M_i &= \Big[\, \textstyle\sum_{j \neq i} \beta_{ij}^{1/\rho}\, C_{T_j i}^{\frac{\rho-1}{\rho}} \Big]^{\frac{\rho}{\rho-1}}
&& \text{(across origins, elasticity } \rho, \text{ weights } \textstyle\sum_{j\neq i}\beta_{ij}=1)
\label{eq:inner}
\end{align}
with the share-linear CES convention (weights enter expenditure shares linearly).

The two elasticities correspond directly to the two forms of adjustment studied below. The elasticity $\eta$ governs the \textbf{home-versus-foreign margin}: following a tariff, how much expenditure moves between domestic and imported tradables. The elasticity $\rho$ governs the \textbf{foreign-supplier reallocation margin}: conditional on importing, how readily expenditure moves from the tariffed exporter toward untreated foreign suppliers. This distinction is irrelevant in a two-country model, because there is only one foreign supplier. It becomes economically active as soon as a third country is introduced.

With the roles of $\eta$ and $\rho$ fixed, the special cases can be stated. The symmetric benchmark used below imposes more than equal origin weights $\beta_{ij} = 1/(N-1)$. Countries also share the preference parameters $(\alpha_T, \alpha_D, \eta, \rho)$ and have equal labor endowments and productivities, so all countries have equal size at the free-trade baseline. Home and foreign tradables remain asymmetric through the home weight $\alpha_D$. Equal treatment of all tradable origins would be a still stronger restriction, since it additionally requires $\rho = \eta$ and $\alpha_D = 1/N$; it is not imposed. The single-elasticity CES is the special case $\rho = \eta$. Cobb--Douglas substitution between the domestic tradable and the import composite is the special case $\eta = 1$; it fixes the elasticity, not the weight $\alpha_D$.

\subsection{Tariffs, income, trade flows, and closure}

Tariffs $\tau_{ij} \geq 0$ are ad valorem, levied by $i$ on imports from $j$. The consumer price of $T_j$ in $i$ is $p_{ij} = (1+\tau_{ij})\, e_{ij}$, in $i$'s currency, where $e_{ij}$ is defined in Section 2.4. Nothing requires $\tau_{ii} = 0$; a tax on home tradables, $p_{ii} = 1+\tau_{ii}$, rebated like the rest, fits the same structure. We set $\tau_{ii} = 0$ throughout, so $\tau$ is a border instrument and $p_{ii} = 1$.

CES price indices in $i$'s currency, from \eqref{eq:mid}--\eqref{eq:inner}:
\begin{equation}
P_{M_i} = \Big[ \textstyle\sum_{j\neq i} \beta_{ij}\, p_{ij}^{\,1-\rho} \Big]^{\frac{1}{1-\rho}},
\qquad
P_{T_i} = \Big[ \alpha_D\, p_{ii}^{\,1-\eta} + (1-\alpha_D)\, P_{M_i}^{\,1-\eta} \Big]^{\frac{1}{1-\eta}},
\label{eq:indices}
\end{equation}
and the consumer price index from \eqref{eq:outer} ($\kappa_i$ collects share constants):
\begin{equation}
P_i = \kappa_i\, P_{T_i}^{\,\alpha_T}\, P_{N_i}^{\,1-\alpha_T} .
\label{eq:cpi}
\end{equation}
With income $I_i$, the share-linear CES demands give expenditure shares of income (rows sum to $\alpha_T$ over tradables):
\begin{align}
s_{ii} &= \alpha_T\, \alpha_D \left(\frac{p_{ii}}{P_{T_i}}\right)^{1-\eta},
\label{eq:shome}\\
s_{ij} &= \alpha_T\, (1-\alpha_D) \left(\frac{P_{M_i}}{P_{T_i}}\right)^{1-\eta}
\beta_{ij} \left(\frac{p_{ij}}{P_{M_i}}\right)^{1-\rho}, \qquad j \neq i .
\label{eq:sfor}
\end{align}
Shares are tariff-inclusive: $s_{ij} I_i$ is spending at consumer prices; $s_{ij} I_i / (1+\tau_{ij})$ is the producer-price (border) value.

Tariff revenue is rebated lump-sum. A fraction $\tau_{ij}/(1+\tau_{ij})$ of spending $s_{ij} I_i$ is revenue, so
\begin{equation}
I_i = w_i L_i + \sum_{j \neq i} \frac{\tau_{ij}}{1+\tau_{ij}}\, s_{ij} I_i
\quad\Longrightarrow\quad
I_i = \frac{w_i L_i}{1 - \sum_{j\neq i} \dfrac{\tau_{ij}}{1+\tau_{ij}}\, s_{ij}} .
\label{eq:income}
\end{equation}
The shares depend only on prices, so \eqref{eq:income} is a closed form.

Trade balances are measured at producer prices, so tariff revenue is a domestic transfer rather than an international payment. In a common currency,
\begin{equation}
\mathrm{TB}_i = \underbrace{\sum_{k \neq i} f_{ki}}_{\text{exports}}
\;-\; \underbrace{\sum_{j \neq i} f_{ij}}_{\text{imports}},
\qquad
f_{ij} \equiv \frac{s_{ij}\, I_i}{1+\tau_{ij}}
\quad \text{($i$'s imports from $j$, border value)} .
\label{eq:tb}
\end{equation}
By Walras' law $\sum_i \mathrm{TB}_i = 0$, so only $N-1$ conditions are independent, matching the $N-1$ free exchange rates.

\begin{assumption}[Closure]
There are no internationally traded assets: each country's trade balance is zero every period,
$\mathrm{TB}_i = 0$ for all $i$.
\label{as:closure}
\end{assumption}
\noindent Assumption \ref{as:closure} closes the model and pins down the $N-1$ independent relative exchange rates. Equilibrium: given $\{\tau_{ij}\}$, the rates $\{e_{Aj}\}_{j\neq A}$ solve $\mathrm{TB}_i = 0$ for $i \neq A$.

This closure is important for interpreting everything that follows. A tariff first changes trade flows at the existing exchange rates. Exchange rates then move until the resulting trade imbalances disappear.

\subsection{Exchange-rate concepts and sign convention}

$e_{ij}$ = units of $i$'s currency per unit of $j$'s currency, with $e_{ii} = 1$ and no-arbitrage $e_{ij} = e_{ik}e_{kj}$. The sign convention is used throughout: \textbf{a rise in $e_{ij}$ is a depreciation of $i$ against $j$}. Thus a negative response of $e_{Aj}$ to $A$'s tariff is an appreciation of $A$ against $j$. Perturbations are written $\nu_j \equiv d\log e_{Aj}$ with $\nu_A = 0$ ($A$'s currency is the numeraire).

The objects to be signed:
\begin{align}
\omega_{ij} &\equiv \frac{e_{ij} w_j}{w_i} = \frac{A_{T_j}}{A_{T_i}}\, e_{ij} \;\propto\; e_{ij}
&& \text{(relative wage, common currency)}
\label{eq:relwage}\\[2pt]
\mathrm{ToT}_{ij} &\equiv \frac{P_{T_i}}{e_{ij} P_{T_j}} = \frac{1}{e_{ij}}
&& \text{(bilateral terms of trade, pre-tariff prices)}
\label{eq:tot}\\[2pt]
q_{ij} &\equiv \frac{e_{ij} P_j}{P_i}
&& \text{(bilateral real exchange rate)}
\label{eq:rer}\\[2pt]
d\nu_i^{E} &\equiv \sum_{j \neq i} w_{ij}\, d\log e_{ij},
\qquad
w_{ij} = \frac{f_{ij} + f_{ji}}{\sum_{k\neq i} (f_{ik} + f_{ki})}
&& \text{(NEER, trade weights)}
\label{eq:neer}\\[2pt]
d q_i^{E} &\equiv \sum_{j \neq i} w_{ij}\, d\log q_{ij}
&& \text{(REER, same weights)}
\label{eq:reer}
\end{align}
where $f_{ij}$ is the border import value defined in \eqref{eq:tb}. The choice of currency cancels from the weights. Effective rates are depreciation-positive: $d\nu_i^E > 0$ means $i$'s currency loses value on average. In the symmetric model $w_{ij} = 1/(N-1)$. The NEER is defined here for completeness; its economic role is developed in Section 6, after the three-country bilateral ambiguity has created a reason to care about aggregation.

Prices and demands depend on $w_j$ and $e_{ij}$ only through the products $e_{ij}w_j$. The model therefore does not separate relative wage adjustment from nominal exchange-rate adjustment. Fixing $w_i = A_{T_i}$ and letting $e$ adjust, as here, is equivalent to fixing $e$ and letting relative wages adjust. With one factor and constant returns, producer prices are unit labor costs, so \eqref{eq:relwage}--\eqref{eq:tot} give
\begin{equation}
d\log e_{ij} \;=\; d\log \omega_{ij} \;=\; -\,d\log \mathrm{ToT}_{ij} .
\label{eq:onetoone}
\end{equation}
The nominal exchange rate, the common-currency relative wage, and the inverse terms of trade move one-for-one. Every result below about $e_{ij}$ can therefore be read as a statement about any of the three.

\section{Two Countries: The Conventional Appreciation Result}

\subsection{Why a tariff appreciates the currency}

Begin with countries $A$ and $B$, with $A$ imposing a small tariff $d\tau$ on imports from $B$ at a symmetric balanced free-trade baseline (incomes normalized to one). There is only one exchange rate and only one foreign supplier. The distinction between aggregate foreign adjustment and reallocation across foreign countries therefore does not yet exist, and $\rho$ is inoperative. Write
\begin{equation}
m \equiv \alpha_T(1-\alpha_D)
\label{eq:m-def}
\end{equation}
for the baseline import share of income, so $s_{AB} = s_{BA} = m$.

At unchanged exchange rates, the tariff reduces $A$'s imports from $B$ and creates a trade-balance disturbance. Equilibrium requires the exchange rate to move so that trade balance is restored. The linearization makes both steps explicit. With $\nu_B = d\log e_{AB}$, price changes in own currency are $d\log p_{AB} = \nu_B + d\tau$ and $d\log p_{BA} = -\nu_B$; with a single import origin these are also the import-index changes, and $d\log P_{T_i} = (1-\alpha_D)\, d\log P_{M_i}$. The expenditure shares respond only on the home-versus-import margin,
\begin{equation}
d\log s_{AB} = (1-\eta)\,\alpha_D\,(\nu_B + d\tau),
\qquad
d\log s_{BA} = -(1-\eta)\,\alpha_D\,\nu_B ,
\label{eq:2c-shares}
\end{equation}
incomes in the common currency move by $d\log I_A = m\, d\tau$ (the rebated revenue) and $d\log I_B = \nu_B$, and the border value of imports loses the tariff wedge $d\tau$. Differentiating $\mathrm{TB}_A$ in \eqref{eq:tb} flow by flow and collecting terms:
\begin{equation}
d\,\mathrm{TB}_A
= \underbrace{m\, D_2\, \nu_B}_{\text{adjustment}}
\;+\; \underbrace{m\, \rho_1^*\, d\tau}_{\text{impact disturbance}}
= 0,
\qquad
D_2 \equiv 2\alpha_D(\eta-1) + 1 ,
\label{eq:2c-collected}
\end{equation}
which defines the two objects that organize everything below. Setting $\nu_B = 0$ in \eqref{eq:2c-collected} shows that at unchanged exchange rates the tariff improves $A$'s balance by $m\rho_1^* d\tau$; the coefficient $mD_2$ is the rate at which a depreciation of $A$ improves its balance. Solving:
\begin{equation}
\boxed{\;
\frac{d\log e_{AB}}{d\tau}\bigg|_{\tau=0}
= -\,\frac{\rho_1^*}{D_2} .
\;}
\label{eq:2c-slope}
\end{equation}

The numerator is
\begin{equation}
\rho_1^*
= 1 + \alpha_D(\eta-1) - \alpha_T(1-\alpha_D)
= \underbrace{\alpha_D\,\eta}_{\text{home switching}} + \underbrace{(1-\alpha_D)(1-\alpha_T)}_{\text{expenditure absorption}}
\;>\; 0
\label{eq:rho1-identity}
\end{equation}
for all admissible parameters. The decomposition is economically useful. The first term captures substitution toward the home tradable; the second captures the expenditure absorbed by the non-tradable margin. Both imply that the tariff creates a surplus requiring an appreciation of $A$, provided the exchange-rate adjustment mechanism has the conventional sign. The denominator $D_2$ is precisely the Marshall--Lerner margin, as the next subsection shows. The notation $\rho_1^*$ anticipates its later role: in the multilateral model this same object turns out to be the per-partner unit in which every threshold is measured.

\begin{result}[Conventional wisdom, $N=2$]
In the two-country model, the exchange-rate adjustment condition is the single inequality $D_2 > 0$. This holds automatically for $\eta \geq 1$ and fails only under strong complementarity, $\eta < 1 - 1/(2\alpha_D)$. Section 3.2 shows it is exactly the textbook condition $\varepsilon_X + \varepsilon_M > 1$. Under it, a unilateral import tariff appreciates the tariffing country's currency, since $\rho_1^* > 0$ by \eqref{eq:rho1-identity}.
\end{result}

The sign therefore has two separate ingredients. The Marshall--Lerner condition signs the \emph{denominator}: exchange rates move trade balances in the normal direction. The identity $\rho_1^* > 0$ signs the \emph{numerator}: the trade-balance disturbance the tariff creates. At $N = 2$ both are unambiguous, and the conclusion is:

\begin{center}
\textbf{A unilateral import tariff appreciates the tariffing country's currency against its only trading partner.}
\end{center}

\noindent By \eqref{eq:onetoone}, the appreciation is equivalently an improvement in $A$'s terms of trade and a rise in its common-currency relative wage; real exchange rates are treated in Section 9. The natural question is what happens when ``its only trading partner'' becomes ``one of several trading partners.''

\subsection{Connection to the textbook Marshall--Lerner condition}
\label{sec:ml}

This subsection imposes nothing new. It restates the condition $D_2 > 0$ in the classical elasticity form in which the conventional wisdom is usually asserted, and then shows that in the nested model the elasticity condition and the denominator $D_2$ are the same object.

Reduced form: export volume $X(p^*)$ with $p^* = 1/e$, import volume $M(p)$ with $p = (1+\tau)e$, elasticities $\varepsilon_X \equiv -X'p^*/X \geq 0$, $\varepsilon_M \equiv -M'p/M \geq 0$. The balance in $A$'s currency at producer prices:
\begin{equation}
\mathrm{TB}_A(e, \tau) = X(1/e) - e\, M\big((1+\tau)e\big) .
\label{eq:ml-tb}
\end{equation}
Differentiate at $\tau = 0$, balanced trade ($X = eM$):
\begin{equation}
\frac{\partial \mathrm{TB}_A}{\partial e}
= \underbrace{\frac{\varepsilon_X X}{e}}_{\text{export volume}}
+ \underbrace{\varepsilon_M M}_{\text{import volume}}
- \underbrace{M}_{\text{valuation}}
= M \big( \varepsilon_X + \varepsilon_M - 1 \big) .
\label{eq:ml}
\end{equation}
A depreciation improves the balance iff $\varepsilon_X + \varepsilon_M > 1$: the volume responses must offset the one-for-one valuation loss. The tariff response at fixed $e$:
\begin{equation}
\frac{\partial \mathrm{TB}_A}{\partial \tau}
= \frac{\varepsilon_M}{1+\tau}\, e M > 0
\quad\text{whenever } \varepsilon_M > 0 .
\label{eq:ml-tau}
\end{equation}
Equilibrium $\mathrm{TB}_A(e(\tau), \tau) = 0$ and the implicit function theorem give
\begin{equation}
\frac{de}{d\tau} = -\, \frac{\partial \mathrm{TB}_A / \partial \tau}{\partial \mathrm{TB}_A / \partial e} :
\label{eq:ml-ift}
\end{equation}
the tariff creates a surplus at the initial rate \eqref{eq:ml-tau}; the rate must move in the direction that worsens the balance; under Marshall--Lerner that direction is appreciation. The condition is necessary and sufficient because $\operatorname{sign}(de/d\tau) = -\operatorname{sign}(\varepsilon_X + \varepsilon_M - 1)$.

In the nested model the condition is not free-standing. Comparing \eqref{eq:ml} with the collected balance \eqref{eq:2c-collected} (per unit of import value $m$):
\begin{equation}
\varepsilon_X + \varepsilon_M - 1 = D_2 = 2\alpha_D(\eta-1) + 1,
\qquad
\varepsilon_X = \varepsilon_M = \alpha_D(\eta - 1) + 1 :
\label{eq:ml-D2}
\end{equation}
the denominator $D_2$ \emph{is} the Marshall--Lerner margin, so the adjustment condition and $\varepsilon_X + \varepsilon_M > 1$ are the same condition. At $\eta = 1$ (Cobb--Douglas), $\varepsilon_X = \varepsilon_M = 1$ and the margin is exactly $1$.

\subsection{Exact Cobb--Douglas benchmark}

At $\eta = 1$ shares are constant: $s_{AB} = s_{BA} = m$. From \eqref{eq:income}, $I_A = w_A L_A / \big[1 - \tfrac{\tau}{1+\tau} m\big]$, and $\mathrm{TB}_A = 0$ in \eqref{eq:tb} reads $e\, m\, w_B L_B = m I_A/(1+\tau)$:
\begin{equation}
\boxed{\;
e(\tau) = \frac{w_A L_A}{w_B L_B}\cdot \frac{1}{1 + (1-m)\tau}
\;}
\qquad
\frac{d\log e}{d\tau}\bigg|_{\tau=0} = -(1-m) = -\frac{\rho_1^*}{D_2}\bigg|_{\eta=1} ,
\label{eq:2c-exact}
\end{equation}
strictly decreasing in $\tau$ globally. The linearized result \eqref{eq:2c-slope} is therefore exact in sign at all tariff levels in the Cobb--Douglas case. Note that $\eta = 1$ pins the expenditure shares at the weights $(\alpha_D, 1-\alpha_D)$ but does not restrict the weights themselves. Cobb--Douglas means unit elasticity, not equal shares; equal home and import spending, $\alpha_D = 1/2$, is a further special case.

\section{Three Countries: Where Multilateralism First Matters}

\subsection{What changes when we add an untreated country?}

Introduce a third country $C$, symmetric with $B$, but leave $C$ untariffed: $A$ imposes the tariff only on $B$. This apparently small change introduces an economic margin that cannot exist in the two-country model. Expenditure that $A$ no longer directs toward $B$ need not return entirely to $A$'s domestic good. It can instead be redirected toward $C$.

The tariff therefore has two conceptually different effects:
\[
\underbrace{\text{change in $A$'s demand for foreign goods as a whole}}_{\text{home-versus-foreign adjustment}}
\qquad\text{and}\qquad
\underbrace{\text{change in the allocation of foreign demand between $B$ and $C$}}_{\text{foreign-supplier reallocation}} .
\]
The second margin is governed by $\rho$. It is the source of the new bilateral results.

\subsection{Impact effects before exchange rates move}

Before solving for exchange rates, ask what the tariff does to trade balances at the initial rates. Holding exchange rates fixed, $B$ necessarily loses: its exports to $A$ fall, so its impact disturbance is negative. For $C$ the effects oppose each other. $C$ gains because some of $A$'s expenditure is diverted away from $B$ toward $C$, a force that strengthens with $\rho$. But $C$ also loses because $B$'s income falls and $B$ reduces its demand for $C$'s goods, and because part of $A$'s expenditure switches toward $A$'s own domestic tradable rather than toward $C$. The impact effect on $C$ is therefore ambiguous. The linearization confirms this and locates the crossover exactly at $\rho_1^*$.

Set $N = 3$, symmetric weights $\beta_{ij} = 1/2$, equal sizes, and perturb by $d\tau = dt_{AB}$. Baseline shares are $s_{ii} = \alpha_T\alpha_D$ and $s_{ij} = m/2$. Log price changes in own currency are $dp_{ij} = \nu_j - \nu_i + dt_{ij}$ with $\nu_A = 0$, so the import price indices move by
\begin{equation}
d\log P_{M_A} = \tfrac{1}{2}(\nu_B + \nu_C + d\tau), \qquad
d\log P_{M_B} = \tfrac{1}{2}\nu_C - \nu_B, \qquad
d\log P_{M_C} = \tfrac{1}{2}\nu_B - \nu_C ,
\label{eq:3c-indices}
\end{equation}
and incomes in the common currency by $d\log I_A = \tfrac{m}{2} d\tau$, $d\log I_B = \nu_B$, $d\log I_C = \nu_C$. Substituting into the differentiated trade balances for $i = B, C$ (shares respond on the cross-origin margin with elasticity $\rho$ and on the home-versus-import margin with elasticity $\eta$, exactly as in \eqref{eq:2c-shares}) and collecting terms in $(\nu_B, \nu_C, d\tau)$ yields the linear system
\begin{equation}
J
\begin{pmatrix} \nu_B \\ \nu_C \end{pmatrix}
= -F\, d\tau,
\qquad
J = -\frac{mD}{4}\begin{pmatrix} 2 & -1 \\ -1 & 2 \end{pmatrix},
\qquad
F = -\frac{m}{4}\begin{pmatrix} \rho + \rho_1^* \\ \rho_1^* - \rho \end{pmatrix},
\label{eq:3c-JF}
\end{equation}
with $D \equiv D_3 = 3\alpha_D(\eta-1) + \rho + 1$. Here $J$ collects the exchange-rate coefficients and $F$ the tariff coefficients at unchanged exchange rates; these objects are generalized in Section 5. The elasticities enter $J$ only through the scalar $D$; the matrix part is pure counting.

The vector $F$ is the formal statement of the impact effects:
\begin{equation}
F_B = -\frac{m}{4}\,(\rho + \rho_1^*) < 0,
\qquad
F_C = \frac{m}{4}\,(\rho - \rho_1^*) .
\label{eq:3c-F}
\end{equation}
$B$'s balance necessarily worsens. $C$'s balance improves exactly when $\rho > \rho_1^*$: sufficiently strong substitutability among foreign suppliers makes the untreated country a beneficiary of the tariff at unchanged exchange rates. But a positive impact effect on $C$ is not yet sufficient for $A$ to depreciate against $C$. Exchange rates jointly restore all countries' balances, so $C$'s equilibrium response also reflects its exposure to $B$'s deterioration.

\subsection{Bilateral equilibrium responses}

Inverting \eqref{eq:3c-JF} with $\left(\begin{smallmatrix} 2 & -1 \\ -1 & 2\end{smallmatrix}\right)^{-1} = \tfrac{1}{3}\left(\begin{smallmatrix} 2 & 1 \\ 1 & 2 \end{smallmatrix}\right)$:
\begin{equation}
\begin{pmatrix} \nu_B \\ \nu_C \end{pmatrix}
= (-J)^{-1} F\, d\tau
= -\frac{1}{3D}
\begin{pmatrix} 2 & 1 \\ 1 & 2 \end{pmatrix}
\begin{pmatrix} \rho + \rho_1^* \\ \rho_1^* - \rho \end{pmatrix} d\tau
= \frac{1}{3D}
\begin{pmatrix} -(\rho + 3\rho_1^*) \\ \rho - 3\rho_1^* \end{pmatrix} d\tau .
\label{eq:3c-solution}
\end{equation}

\begin{result}[Threshold]
At the symmetric point the exchange-rate adjustment condition reduces to $D > 0$, automatic for $\eta \geq 1$ and $\rho > 0$. Under it,
\begin{equation}
\frac{d\log e_{AB}}{d\tau} = -\,\frac{\rho + 3\rho_1^*}{3D} < 0,
\qquad
\boxed{\;\frac{d\log e_{AC}}{d\tau} = \frac{\rho - 3\rho_1^*}{3D}\;}
\qquad
\frac{d\log e_{BC}}{d\tau} = \frac{2\rho}{3D} > 0 .
\label{eq:3c-slopes}
\end{equation}
$A$ appreciates against the target $B$ for all parameters; $A$ depreciates against the untariffed $C$ iff $\rho > \rho_3^* \equiv 3\rho_1^*$; $B$ depreciates against $C$ for all parameters.
\end{result}

Three results emerge. First, $A$ always appreciates against the tariff target $B$: the familiar bilateral result survives for the country directly affected by the tariff. Second, $B$ always depreciates against the untreated country $C$: the tariff lowers demand for $B$ relative to $C$, and equilibrium requires a corresponding deterioration in $B$'s relative price. Third, and most importantly, $A$'s exchange rate against $C$ is ambiguous. When cross-origin substitution is sufficiently strong, trade diversion toward $C$ is large enough that $C$ appreciates not only against $B$ but also against $A$. The country imposing the tariff then \textbf{appreciates against its target while depreciating against an untreated trading partner}. This is the precise point at which the literal two-country conventional wisdom fails.

\subsection{The meaning of the bilateral threshold}

The threshold $3\rho_1^*$ is the point at which foreign-supplier reallocation becomes strong enough to dominate the forces working against $C$. The structure of the solution shows why the threshold sits at three times the impact crossover. Writing the equilibrium response through the nonnegative weights of $(-J)^{-1}$,
\begin{equation}
\frac{d\nu_C}{d\tau} \propto F_B + 2F_C \propto \rho - 3\rho_1^* :
\label{eq:3c-signrule}
\end{equation}
$C$'s equilibrium response mixes its own impact gain with $B$'s impact loss. $F_C > 0$ at $\rho > \rho_1^*$, but reversal of $e_{AC}$ requires the diversion gain to outweigh $C$'s exposure to $B$'s deterioration, which happens only at $\rho > 3\rho_1^*$. A higher $\rho$ makes foreign suppliers easier to substitute across, so more of the expenditure displaced from the tariffed $B$ is redirected toward $C$; the response against $C$ changes sign only when that diversion becomes strong enough.

This is not a failure of Marshall--Lerner logic. The ambiguity arises because the tariff creates a mixed pattern of trade-balance disturbances across countries. That observation motivates a general way to organize the comparative statics, introduced next.

\section{Disturbance and Adjustment: The General Multilateral Framework}

The three-country case suggests a general decomposition. A tariff first creates a vector of trade-balance disturbances at unchanged exchange rates. Exchange rates then move to eliminate those disturbances. The comparative static can be read as
\[
\text{tariff}
\;\longrightarrow\;
\text{trade-balance disturbances } F
\;\longrightarrow\;
\text{exchange-rate adjustment } (-J)^{-1}F .
\]
This section derives that structure for general $N$ and arbitrary weights, states the condition under which the adjustment mechanism is regular, and gives the resulting sign rule. The systems of Sections 3 and 4 are its $N = 2$ and $N = 3$ cases.

\subsection{General linearization}

Evaluate at $\tau_{ij} = 0$ with balanced trade, and normalize baseline incomes. The baseline objects are the home share within tradables $h_i = \alpha_D$, the within-import shares $b_{ij} = \beta_{ij}$, the income shares $s_{ij} = \alpha_T(1-\alpha_D) b_{ij} = m\, b_{ij}$, and the flow values $f_{ij} = s_{ij} I_i$ from \eqref{eq:tb} at $\tau = 0$. A country-specific home share $h_i$ carries through the derivation unchanged. Balance requires $\sum_k f_{ki} = \sum_j f_{ij}$.

Perturb by tariffs $dt_{ij}$ and collect exchange-rate changes $\nu_j = d\log e_{Aj}$. Log price changes in $i$'s own currency ($dp_{ii} = 0$):
\begin{equation}
dp_{ij} = \nu_j - \nu_i + dt_{ij} .
\label{eq:dp}
\end{equation}
Index changes, from \eqref{eq:indices}:
\begin{equation}
d\log P_{M_i} = \sum_{j\neq i} b_{ij}\, dp_{ij},
\qquad
d\log P_{T_i} = (1-\alpha_D)\, d\log P_{M_i} .
\label{eq:dindex}
\end{equation}
Share changes, from \eqref{eq:shome}--\eqref{eq:sfor}, using $d\log P_{M_i} - d\log P_{T_i} = \alpha_D\, d\log P_{M_i}$:
\begin{align}
d\log s_{ii} &= -(1-\eta)(1-\alpha_D)\, d\log P_{M_i},
\label{eq:dshome}\\
d\log s_{ij} &= \underbrace{(1-\rho)\big( dp_{ij} - d\log P_{M_i} \big)}_{\text{across origins}}
+ \underbrace{(1-\eta)\,\alpha_D\, d\log P_{M_i}}_{\text{home vs.\ imports}} .
\label{eq:dsfor}
\end{align}
Income changes in the common currency, from \eqref{eq:income} (revenue is first order only for tariffing countries):
\begin{equation}
d\log I_i = \nu_i + \sum_{j\neq i} s_{ij}\, dt_{ij} .
\label{eq:dincome}
\end{equation}
Differentiating \eqref{eq:tb}, each flow contributes value $\times$ log change:
\begin{equation}
d\,\mathrm{TB}_i
= \sum_{k \neq i} f_{ki} \big( d\log s_{ki} + d\log I_k - dt_{ki} \big)
- \sum_{j \neq i} f_{ij} \big( d\log s_{ij} + d\log I_i - dt_{ij} \big) .
\label{eq:dtb}
\end{equation}
Equations \eqref{eq:dp}--\eqref{eq:dtb} are linear in the exchange-rate changes and the tariff changes. Stack the free rates in the vector $d\nu \equiv (\nu_j)_{j \neq A} \in \mathbb{R}^{N-1}$ and all directed bilateral tariff changes in $d\boldsymbol{\tau} \equiv (dt_{jk})_{j \neq k}$. Reading the coefficients off \eqref{eq:dtb}, the linearized equilibrium conditions for $i \neq A$ are
\begin{equation}
J\, d\nu + G\, d\boldsymbol{\tau} = 0,
\qquad
J_{il} \equiv \frac{\partial\, d\mathrm{TB}_i}{\partial \nu_l},
\qquad
G_{i,(jk)} \equiv \frac{\partial\, d\mathrm{TB}_i}{\partial\, dt_{jk}}\bigg|_{\nu \text{ fixed}} .
\label{eq:JFdef}
\end{equation}
$J$ is the $(N-1)\times(N-1)$ \textbf{adjustment mechanism}: how trade balances respond to exchange rates. $G$ maps any pattern of tariff changes into the trade-balance disturbances it creates before exchange rates move. A policy experiment is a direction $a = (a_{jk})$ scaled by a single number, $d\boldsymbol{\tau} = a\, d\tau$. The unilateral tariff of the preceding sections is $a_{AB} = 1$ with all other entries zero; the trade war of Section 8 is $a_{AB} = a_{BA} = 1$. The \textbf{impact disturbance} of an experiment is
\begin{equation}
F \equiv G a,
\qquad
F_i = \frac{\partial\, d\mathrm{TB}_i}{\partial\, d\tau}\bigg|_{\nu = 0} .
\label{eq:F-def}
\end{equation}
One sign convention is fixed and used throughout: $F$ is the impact effect of the experiment on the trade balances at unchanged exchange rates. Equation \eqref{eq:JFdef} then reads
\begin{equation}
J\, d\nu = -F\, d\tau
\quad\Longrightarrow\quad
\frac{d\nu}{d\tau} = (-J)^{-1} F = \frac{\operatorname{adj}(-J)\, F}{\det(-J)},
\label{eq:system}
\end{equation}
where $\operatorname{adj}(-J)$ is the adjugate (transpose of the cofactor matrix). Different tariff configurations are different directions $a$ applied to the same $G$. This is the superposition used in Section 8. The entries of $J$ and $G$ are linear in $\eta$ and $\rho$, with baseline shares as coefficients. At $N = 3$ with symmetric weights, \eqref{eq:system} is exactly \eqref{eq:3c-JF}; at $N = 2$ it is the scalar equation \eqref{eq:2c-collected}.

\subsection{The multilateral Marshall--Lerner condition and the sign rule}

For the familiar comparative statics to operate normally, depreciation must improve the trade balance of the depreciating country. In a two-country model this is the scalar Marshall--Lerner condition $D_2 > 0$. With several exchange rates, the corresponding condition applies to the whole adjustment system: we require $-J$ to be a nonsingular M-matrix,
\begin{equation}
\det(-J) > 0 \ \text{(all principal minors positive)},
\qquad
(-J)^{-1} \geq 0 \ \text{elementwise},
\qquad
\operatorname{Re}\lambda(J) < 0 ,
\label{eq:invpos}
\end{equation}
three equivalent statements that we refer to jointly as the \emph{multilateral Marshall--Lerner condition}. Section 5.3 derives \eqref{eq:invpos} from primitives and shows that regularity, local stability, and inverse positivity are one notion, not three. Here we record its economic implication.

\begin{result}[Sign rule]
Under the multilateral Marshall--Lerner condition \eqref{eq:invpos},
\begin{equation}
\operatorname{sign}\!\left( \frac{d\nu_i}{d\tau} \right)
= \operatorname{sign}\!\Big( \sum_j \omega_{ij} F_j \Big),
\qquad \omega_{ij} \equiv \big[(-J)^{-1}\big]_{ij} \geq 0 .
\label{eq:signrule}
\end{equation}
Every response is a nonnegative combination of the impact effects. If all $F_j$ share one sign, every exchange rate moves the same way; opposite-signed responses require an impact vector with mixed signs.
\end{result}

This clarifies the source of bilateral ambiguity. It does not come from a perverse adjustment matrix. It comes from the direction of $F$. For the discriminatory tariff, the impact vector is generically mixed, exactly as seen concretely in Section 4.2: $F_B < 0$ always, while $F_C$ combines a diversion gain (increasing in $\rho$) with reduced sales to a poorer $B$, giving $F_C \propto \rho - \rho_1^*$ in the symmetric three-country case. The thresholds derived in these notes are the points where a weighted impact sum $\sum_j \omega_{ij} F_j$ crosses zero; the three-country threshold \eqref{eq:3c-signrule} is the explicit example. Because $F$ is linear in $\rho$ and $\operatorname{adj}(-J)$ has degree $N-2$ in $\rho$, each weighted sum is a polynomial of degree $N-1$ in $\rho$. It has one relevant root at $N = 3$ and can change sign more than once at $N \geq 4$.

\subsection{Technical foundations: from primitives to the M-matrix property}

This subsection derives \eqref{eq:invpos} from the model's structure plus one assumption. Write $\widetilde J$ for the full $N \times N$ Jacobian, $\widetilde J_{il} = \partial\, d\mathrm{TB}_i / \partial \nu_l$ for all $i, l \in \{1, \dots, N\}$, before dropping country $A$. The $J$ of \eqref{eq:JFdef} is $\widetilde J$ with row and column $A$ deleted. Two properties of $\widetilde J$ are identities of the model, not assumptions:
\begin{align}
\sum_{l} \widetilde J_{il} &= 0 \quad \forall i
&& \text{(homogeneity: \eqref{eq:dp} depends only on differences } \nu_j - \nu_i)
\label{eq:hom}\\
\sum_{i} \widetilde J_{il} &= 0 \quad \forall l
&& \text{(Walras' law: } \textstyle\sum_i \mathrm{TB}_i \equiv 0) .
\label{eq:walras-col}
\end{align}

\begin{assumption}[Gross substitutability]
$\widetilde J_{il} \geq 0$ for all $l \neq i$: a depreciation of any country $l$ (a rise in $\nu_l$) does not worsen any other country $i$'s trade balance.
\label{as:gs}
\end{assumption}
\noindent The condition applies to every off-diagonal entry, for all country pairs, not only $A$'s partners. In the explicit symmetric systems of Sections 4 and 7, every off-diagonal entry is a positive multiple of $D_3$ or $D_N$. Assumption \ref{as:gs} there is exactly $D_N > 0$, which holds for all $\eta \geq 1$.

Combining \eqref{eq:hom} with Assumption \ref{as:gs} gives the diagonal, and then the reduced system inherits a dominant diagonal:
\begin{equation}
\widetilde J_{ii} = -\sum_{l \neq i} \widetilde J_{il} \;\leq\; 0,
\qquad\qquad
|J_{ii}| = \sum_{l \neq i,\,A} J_{il} \;+\; \widetilde J_{iA}
\;\geq\; \sum_{l \neq i,\,A} J_{il} .
\label{eq:diagdom}
\end{equation}
The first equation is the set of bilateral Marshall--Lerner conditions, one per country: own depreciation improves own balance. The second equation says that each row of $-J$ dominates its off-diagonal sum. The dominance is strict in every row $i$ with $\widetilde J_{iA} > 0$, that is, for every country whose balance responds to $A$'s exchange rate. Diagonal dominance is therefore neither an independent assumption nor a restatement of the identities. It follows from homogeneity \eqref{eq:hom} together with gross substitutability, once the numeraire column is dropped. Walras' law \eqref{eq:walras-col} delivers the same dominance columnwise.

\begin{result}[Multilateral Marshall--Lerner condition from primitives]
Under Assumption \ref{as:gs}, if $J$ is irreducible (a connected trade network) with $\widetilde J_{iA} > 0$ for some $i$, then $-J$ is a nonsingular M-matrix, and the three equivalent conditions \eqref{eq:invpos} hold.
\label{res:mml}
\end{result}
\noindent The proof is the standard characterization of irreducibly diagonally dominant Z-matrices; \eqref{eq:diagdom} supplies the dominance. Three clarifications of the logic:
\begin{itemize}
\item Gross substitutability is a \emph{sufficient primitive}, not the condition itself. Assumption \ref{as:gs} implies \eqref{eq:invpos}, but \eqref{eq:invpos} can hold without it. Everything in these notes uses only \eqref{eq:invpos}.
\item The eigenvalue statement in \eqref{eq:invpos} is local stability of the \emph{tatonnement}, the price-adjustment process $\dot\nu_i = k_i\, \mathrm{TB}_i$ with any speeds $k_i > 0$, under which a country in surplus sees its currency appreciate. Stability is therefore not a requirement layered on top of \eqref{eq:invpos}. It is one of its equivalent faces, and M-matrices are stable for every choice of speeds.
\item At $N = 2$, all of \eqref{eq:invpos} collapses to the single inequality $J < 0$, that is, $D_2 > 0$. Its classical elasticity form $\varepsilon_X + \varepsilon_M > 1$ was derived in Section 3.2.
\end{itemize}

\section{Aggregation: Why the Conventional Result Survives in the NEER}

The three-country bilateral results create a natural question. If $A$ appreciates against $B$ but can depreciate against $C$, what happens to the overall value of $A$'s currency? With symmetric trade weights, $A$'s nominal effective exchange-rate response is the average of its bilateral responses. Decompose the solution \eqref{eq:3c-solution} into an aggregate and a relative component:
\begin{equation}
\underbrace{\frac{\nu_B + \nu_C}{2}}_{\text{NEER of } A}
= -\,\frac{\rho_1^*}{D}\, d\tau,
\qquad
\underbrace{\nu_B - \nu_C}_{\text{relative}}
= -\,\frac{2\rho}{3D}\, d\tau .
\label{eq:3c-decomp}
\end{equation}
These expressions reveal the structure of the result. The first is the \textbf{aggregate component}: the average value of $A$'s currency against the rest of the world. The second is the \textbf{relative component}: how that adjustment is distributed between $B$ and $C$. Cross-origin substitution $\rho$ directly drives the relative component. As $B$ and $C$ become closer substitutes, the tariff produces progressively more appreciation of $C$ relative to $B$, and eventually this reallocation reverses $A$'s bilateral rate against $C$. But the same $\rho$-terms enter the two bilateral responses with opposite signs and cancel when aggregated. The denominator $D$ still contains $\rho$, so cross-origin substitution affects the magnitude of the NEER response but not its sign. Economically, $\rho$ governs how the exchange-rate adjustment is distributed across foreign partners.

\begin{result}[NEER preservation]
With symmetric weights, $d\nu_A^E/d\tau = -\rho_1^*/D < 0$ for all $\rho$. A unilateral tariff therefore always appreciates the tariffing country's effective exchange rate, even when $\rho > \rho_3^*$ reverses the bilateral rate against $C$. The conventional wisdom fails bilaterally but survives in aggregate, and by \eqref{eq:rho1-identity} the aggregate sign is unconditional.
\end{result}

This gives a different interpretation of the familiar two-country result. In a two-country world there is no distinction between a bilateral exchange rate and an effective exchange rate, so the conventional appreciation result appears bilateral. Once additional trading partners are introduced, the bilateral and aggregate concepts separate. The bilateral result can fail because foreign-supplier reallocation affects different partners differently, while the aggregate result survives because that reallocation cancels across partners.

\begin{center}
\textbf{The two-country appreciation result is best understood as an aggregate home-versus-foreign result, not as a prediction that the tariffing country must appreciate against every foreign currency.}
\end{center}

\paragraph{Where the conditions sit empirically.} Given \eqref{eq:invpos}, $D > 0$ and $\rho_1^* > 0$ hold at every admissible calibration. The signs of $e_{AB}$, $e_{BC}$, and the NEER are therefore parameter-free. The one genuinely ambiguous object is the bilateral rate against the bystander. At the benchmark calibration ($\alpha_D = 0.8$, $\alpha_T = 0.4$, $\eta = 1.5$), $\rho_3^* = 3.96$ sits inside the range of estimated cross-origin trade elasticities, roughly $2$--$8$, with macro estimates at the low end and trade-level estimates at the high end. Either sign of $d\log e_{AC}$ is empirically plausible. By contrast, the trade-war threshold of Section 8, $\rho_1^* = 1.32$, lies beneath essentially all estimates, so the war predictions are clear-cut.

\section{General $N$: Dilution and Aggregation}

The three-country model contains the essential economics. Generalizing to $N$ countries asks how the same mechanisms change as the number of potential alternative suppliers grows. Suppose $A$ tariffs $B$ and the remaining $N-2$ countries are symmetric bystanders; by symmetry all bystanders respond identically, so the system again reduces to two unknowns. Before deriving the result, the economics suggests a prediction. The tariff still diverts expenditure away from $B$, but that diversion is now divided among $N-2$ alternative foreign suppliers. Any individual bystander should receive less of the diversion gain as $N$ rises, so bilateral reversal should become progressively harder, even though the aggregate foreign reallocation need not disappear.

\subsection{The reduced $N$-country system}

Symmetric weights $b_{ij} = 1/(N-1)$, equal sizes, $d\tau = dt_{AB}$; the $N-2$ bystanders are identical with common $\nu_C$. From \eqref{eq:dp}--\eqref{eq:dindex}:
\begin{equation}
d\log P_{M_A} = \frac{\nu_B + d\tau + (N-2)\nu_C}{N-1},
\qquad
d\log P_{M_B} = \frac{(N-2)\nu_C}{N-1} - \nu_B,
\qquad
d\log P_{M_C} = \frac{\nu_B - 2\nu_C}{N-1},
\label{eq:N-indices}
\end{equation}
and $d\log I_A = \tfrac{m}{N-1}d\tau$, $d\log I_B = \nu_B$, $d\log I_C = \nu_C$. Substituting into \eqref{eq:dtb} for $B$ and a representative bystander and collecting terms (each row scaled by $(N-1)/m$; verified symbolically):
\begin{equation}
D_N
\begin{pmatrix} -(N-1) & N-2 \\ 1 & -2 \end{pmatrix}
\begin{pmatrix} \nu_B \\ \nu_C \end{pmatrix}
=
\begin{pmatrix} (N-2)\rho + \rho_1^* \\ \rho_1^* - \rho \end{pmatrix} d\tau ,
\qquad
D_N = N\alpha_D(\eta-1) + (N-2)\rho + 1 .
\label{eq:N-system}
\end{equation}
As at $N=3$, the elasticities enter the adjustment matrix only through the scalar $D_N$; the matrix is pure counting, with determinant $N$. Inverting:
\begin{align}
\frac{d\log e_{AB}}{d\tau}
&= -\,\frac{N\rho_1^* + (N-2)\rho}{N\, D_N} \;<\; 0,
\label{eq:N-eAB}\\[3pt]
\frac{d\log e_{AC}}{d\tau}
&= \frac{\rho - N\rho_1^*}{N\, D_N},
\label{eq:N-eAC}\\[3pt]
\frac{d\log e_{BC}}{d\tau}
&= \frac{(N-1)\,\rho}{N\, D_N} \;>\; 0 .
\label{eq:N-eBC}
\end{align}
At $N = 3$ these are \eqref{eq:3c-slopes}; at $N = 2$, \eqref{eq:N-eAB} is \eqref{eq:2c-slope}, with $\rho$ absent and $D_2 = 2\alpha_D(\eta-1)+1$.

\subsection{The bilateral threshold as dilution}

\begin{result}[$N$-country threshold]
$A$ depreciates against each untariffed bystander iff
\begin{equation}
\boxed{\;\rho > \rho_N^* = N \rho_1^* \;}
\end{equation}
The threshold is linear in $N$ because diverted expenditure splits across $N-2$ bystanders, each receiving weight $1/(N-1)$, so a larger $\rho$ is needed to reverse any one bilateral rate. This is a statement about how the adjustment is distributed across partners, not about the aggregate value of $A$'s currency.
\end{result}

The reason is dilution. With one bystander, expenditure diverted away from $B$ has only one alternative foreign destination. With $N-2$ bystanders, the same foreign reallocation is spread across many suppliers. Each individual bystander receives a smaller diversion gain, and stronger cross-origin substitutability is required to generate enough appreciation to reverse its bilateral rate against $A$.

\subsection{Effective exchange-rate aggregation}

Averaging with symmetric weights:
\begin{equation}
\frac{d\nu_A^E}{d\tau}
= \frac{\nu_B + (N-2)\nu_C}{N-1}
= \frac{-\big[N\rho_1^* + (N-2)\rho\big] + (N-2)\big[\rho - N\rho_1^*\big]}{(N-1)\, N\, D_N}
= \boxed{\; -\,\frac{\rho_1^*}{D_N} \;}
\label{eq:N-neer}
\end{equation}
The $\rho$ terms cancel exactly, and $N$ affects only the magnitude through $D_N$. The cancellation observed with three countries is therefore not accidental; it holds for every $N \geq 2$. The companion relative component is pure $\rho$:
\begin{equation}
\nu_B - \nu_C = -\,\frac{(N-1)\rho}{N D_N}\, d\tau .
\label{eq:N-decomp}
\end{equation}

\begin{result}[Aggregation]
In the symmetric model, for every $N \geq 2$ and every $\rho$, a unilateral tariff appreciates the tariffing country's NEER: $d\nu_A^E/d\tau = -\rho_1^*/D_N < 0$ by \eqref{eq:rho1-identity}. Cross-origin substitution reallocates the adjustment across partners \eqref{eq:N-decomp} and can reverse individual bilateral rates \eqref{eq:N-eAC}, but under symmetry it cannot reverse the NEER. This is the precise sense in which the conventional two-country appreciation result survives multilateralization.
\end{result}

The general model therefore sharpens the decomposition. Home-versus-foreign adjustment maps into the effective exchange rate; reallocation across foreign suppliers maps into the dispersion of bilateral exchange rates. Cross-origin substitution affects the magnitude of the aggregate response through equilibrium adjustment, but under symmetry it cannot change its sign. Its qualitatively new role is to redistribute the adjustment across foreign partners.

\subsection{Single-elasticity restriction: an impossibility result}

The bilateral reversal requires the foreign-supplier reallocation margin to be strong independently of the home-versus-foreign margin. A single-elasticity specification ties the two together. Impose $\rho = \eta = \sigma$. At $N = 3$,
\begin{equation}
\rho_3^* - \sigma
= 3\rho_1^*\big|_{\eta=\sigma} - \sigma
= \sigma(3\alpha_D - 1) + 3(1-\alpha_D)(1-\alpha_T) .
\label{eq:3c-gap}
\end{equation}

\begin{result}[Impossibility]
If $\alpha_D \geq 1/3$, then $\rho_3^* > \sigma$ for every $\sigma > 0$. With one elasticity, the reversal against $C$ is unreachable. The bound involves only $\alpha_D$; no value of $\alpha_T$ is required. As an illustration, suppose expenditure is also equal across the four consumption categories, the three tradables and the non-tradable, so that $\alpha_D = 1/3$ and $\alpha_T = 3/4$. This is the paper's original calibration, not a requirement of the result. Then $\rho_3^* - \sigma = \tfrac{1}{2}$, $D = 2\sigma$, and
\begin{equation}
\frac{d\log e_{AC}}{d\tau} = \frac{\sigma - \rho_3^*}{3D} = -\,\frac{1}{12\sigma} \longrightarrow 0^- .
\end{equation}
\end{result}

At general $N$ the same restriction gives
\begin{equation}
\rho_N^* - \sigma = \sigma(N\alpha_D - 1) + N(1-\alpha_D)(1-\alpha_T)
\;>\; 0 \quad \text{whenever } \alpha_D \geq 1/N .
\label{eq:N-gap}
\end{equation}
The home-bias bound weakens with $N$. Any calibration with home share at least $1/N$ keeps the bilateral reversal invisible to single-elasticity models: when one parameter governs both margins, the cross-origin margin cannot become independently strong enough to produce the reversal over the home-biased region.

\subsection{Correspondence across $N$}

\begin{center}
\begin{tabular}{p{3.6cm}p{3.0cm}p{3.4cm}p{4.2cm}}
\toprule
object & $N = 2$ & $N = 3$ & general $N$ \\
\midrule
ML condition \eqref{eq:invpos} & $\varepsilon_X + \varepsilon_M - 1 = D_2 > 0$ & $D_3 > 0$ & $-J$ a nonsingular M-matrix \\[3pt]
sufficient primitive & $\eta \geq 1$ & $\eta \geq 1$ & gross substitutability (Assumption \ref{as:gs}) \\[3pt]
sign rule & one impact, unconditional & weights $\tfrac{1}{3}\left(\begin{smallmatrix} 2 & 1 \\ 1 & 2\end{smallmatrix}\right)$ & $\operatorname{sign}\big(\sum_j \omega_{ij}F_j\big)$, $\omega \geq 0$ \\[3pt]
bilateral threshold & none (no bystander) & $3\rho_1^*$ & $N\rho_1^*$ (symmetric); degree-$(N-1)$ roots in general \\[3pt]
real bilateral threshold & none & $\rho_q^*(3)$ \eqref{eq:3c-rhoq} & $\rho_q^*(N)$ \eqref{eq:N-rhoq} \\[3pt]
NEER (unilateral) & $-\rho_1^*/D_2$ & $-\rho_1^*/D_3$ & $-\rho_1^*/D_N$: sign invariant \\[3pt]
war threshold & --- & $\rho_1^*$ & $\rho_1^*$, independent of $N$ \\
\bottomrule
\end{tabular}
\end{center}

\section{What Depends on the Tariff Configuration?}

The preceding results concern a discriminatory unilateral tariff. Different tariff configurations activate the two margins differently. Responses are linear in the tariff vector, so any configuration is a sum of directed bilateral shocks, with mirror shocks obtained by relabeling \eqref{eq:3c-solution}; in the notation of Section 5, different configurations are different directions $a$ applied to the same $G$.

\subsection{Broad unilateral protection}

Suppose $A$ raises tariffs uniformly against every foreign supplier. All origins are treated identically, so there is no tariff-induced reallocation across foreign countries: the cross-origin margin is inactive. At $N = 3$, superposition of the tariff on $B$ and its mirror on $C$ gives
\begin{equation}
\frac{d\log e_{Aj}}{d\tau} = \nu_B + \nu_C = -\,\frac{2\rho_1^*}{D} < 0, \qquad j = B, C,
\label{eq:3c-broad}
\end{equation}
and at general $N$, superposition over all $N-1$ partners gives
\begin{equation}
\frac{d\log e_{Aj}}{d\tau} = -\,\frac{(N-1)\rho_1^*}{D_N} \quad \forall j
\;=\; \frac{d\nu_A^E}{d\tau} .
\label{eq:N-broad}
\end{equation}
In both, $\rho$ drops out of the numerator entirely. Every bilateral exchange rate appreciates by the same amount, and the bilateral and effective responses coincide. The comparison with the discriminatory tariff is instructive: bilateral ambiguity is not generated by protection per se. It is generated by \textbf{discrimination across foreign suppliers}.

\subsection{Reciprocal trade war}

Now suppose $B$ retaliates with an equal tariff on $A$. At $N = 3$, adding the mirror shock:
\begin{equation}
\frac{d\log e_{AB}}{d\tau}\bigg|^{w} = 0,
\qquad
\frac{d\log e_{AC}}{d\tau}\bigg|^{w} = \frac{\rho - \rho_1^*}{D},
\label{eq:3c-war}
\end{equation}
and the belligerents' effective rates flip sign at $\rho_1^*$:
\begin{equation}
\frac{d\nu_A^{E}}{d\tau}\bigg|^{w} = \frac{1}{2}\,\frac{\rho - \rho_1^*}{D} > 0
\quad\Longleftrightarrow\quad
\rho > \rho_1^* .
\label{eq:3c-warneer}
\end{equation}
At general $N$:
\begin{equation}
\frac{d\log e_{AB}}{d\tau}\bigg|^{w} = 0,
\qquad
\frac{d\log e_{AC}}{d\tau}\bigg|^{w} = \frac{\rho - \rho_1^*}{D_N},
\label{eq:N-war}
\end{equation}
and the war NEER, $\tfrac{N-2}{N-1}\,\tfrac{\rho-\rho_1^*}{D_N}$, has the sign of $\rho - \rho_1^*$.

\begin{result}[War threshold is independent of $N$]
In a symmetric war, every bystander appreciates against both belligerents, and both belligerents' NEERs depreciate, iff
\begin{equation}
\boxed{\;\rho > \rho_1^* = \frac{\rho_N^*}{N}\;}
\end{equation}
for any $N$. With $e_{AB}$ pinned at zero by symmetry, each bystander's equation decouples, and its sign is the sign of its own impact effect, which does not dilute with $N$. At $\eta = 1$, $\rho_1^* = 1 - m < 1$, so any $\rho \geq 1$ suffices.
\end{result}

\begin{result}[Configuration dependence]
The unilateral tariff appreciates the NEER for all $\rho$; the reciprocal war depreciates both belligerents' NEERs iff $\rho > \rho_1^*$. Retaliation, not protection itself, is what can weaken the tariffing country's currency in aggregate.
\end{result}

Retaliation cancels the direct bilateral component between $A$ and $B$, leaving the foreign-supplier reallocation margin to determine the belligerents' exchange rates against the rest of the world. This also explains why the war threshold is much lower than the unilateral bilateral-reversal threshold $N\rho_1^*$. The latter requires diversion to become strong enough to overcome the direct appreciation against $B$. In a symmetric war that bilateral component is removed by symmetry.

\section{Real Exchange Rates}

The nominal results give the cleanest statement of the mechanism. Real exchange rates add a direct consumer-price effect: the tariff raises the tariffing country's price index even before the nominal exchange rate moves.

\subsection{Bilateral real rates}

From \eqref{eq:rer} and \eqref{eq:cpi}, $d\log q_{Aj} = \nu_j + \alpha_T\big( d\log P_{T_j}^{\mathrm{own}} - d\log P_{T_A}^{\mathrm{own}} \big)$ with own-currency indices \eqref{eq:dindex}. For the unilateral tariff at the symmetric baseline, both bilateral real rates take one form at every $N$ (verified symbolically):
\begin{equation}
d\log q_{Aj} = \Big( 1 - \frac{mN}{N-1} \Big)\, \nu_j \;-\; \frac{m}{N-1}\, d\tau,
\qquad j = B, C .
\label{eq:qgen}
\end{equation}
At $N = 2$ this reads $d\log q_{AB} = (1 - 2m)\, \nu_B - m\, d\tau$ (together with $d\log \omega_{AB} = \nu_B$ and $d\log \mathrm{ToT}_{AB} = -\nu_B$ from \eqref{eq:onetoone}, so at $N = 2$ every definition of ``the exchange rate'' moves together); at $N = 3$,
\begin{equation}
d\log q_{Aj} = \Big(1 - \tfrac{3m}{2}\Big)\, \nu_j \;-\; \tfrac{m}{2}\, d\tau,
\qquad j = B, C .
\label{eq:3c-dq}
\end{equation}
The two pieces of \eqref{eq:qgen} are the pass-through of the nominal rate, attenuated by the tradable exposures, and the direct increase in $A$'s consumer prices, which works toward real appreciation of $A$ at any nominal rate.

The direct price effect shifts the bilateral reversal threshold up. Setting $d\log q_{AC} = 0$ at $N = 3$ with $\nu_C$ from \eqref{eq:3c-solution}: $(2-3m)(\rho - 3\rho_1^*) = 3mD$, linear in $\rho$, so
\begin{equation}
\boxed{\;\rho^*_q = \frac{(2 - 3m)\, 3\rho_1^* + 3m\big[\,3\alpha_D(\eta-1) + 1\,\big]}{2(1 - 3m)}\;}
\qquad (\rho^*_q > \rho_3^* \text{ for } m \in (0, \tfrac{1}{3})) ,
\label{eq:3c-rhoq}
\end{equation}
and at general $N$, using \eqref{eq:N-eAC} (for $m < 1/N$):
\begin{equation}
\boxed{\;
\rho_q^*(N) = \frac{\big[(N-1) - mN\big]\, N\rho_1^* + mN\big[ N\alpha_D(\eta-1) + 1 \big]}{(N-1)(1 - mN)}
\;}
\label{eq:N-rhoq}
\end{equation}
($\rho_q^*(3) = 4.93$, $\rho_q^*(4) = 7.34$, $\rho_q^*(5) = 10.40$ at the benchmark, against nominal thresholds $3.96$, $5.28$, $6.60$). Real reversal requires more diversion than nominal reversal, because part of the nominal movement is absorbed by the tariff's direct effect on $A$'s price index.

\subsection{Effective real rates}

Averaging \eqref{eq:3c-dq} over $j$ gives the REER at $N = 3$:
\begin{equation}
\frac{dq_A^E}{d\tau} = -\Big(1 - \tfrac{3m}{2}\Big)\frac{\rho_1^*}{D} - \frac{m}{2} \;<\; 0
\qquad \text{for all } \rho \ (m < \tfrac{2}{3}) ,
\label{eq:3c-reer}
\end{equation}
and averaging \eqref{eq:qgen} at general $N$:
\begin{equation}
\frac{dq_A^E}{d\tau}
= -\Big( 1 - \frac{mN}{N-1} \Big) \frac{\rho_1^*}{D_N} - \frac{m}{N-1}
\;<\; 0
\qquad \text{for all } \rho, N \ \big(m < \tfrac{N-1}{N}\big) .
\label{eq:N-reer}
\end{equation}
Real effective appreciation is unconditional and strictly stronger than nominal, because the tariff also raises $A$'s consumer prices directly. The effective real appreciation is thus even more robust than the nominal one.

\section{Asymmetry: How Much of the Aggregation Result Survives?}

The exact cancellation underlying the symmetric NEER result relies on a correspondence between bilateral responses and the weights with which they enter the effective exchange rate. Actual trade networks are asymmetric, so the natural question is whether symmetry is essential to the economics or only to the exact closed form. At a general balanced baseline the disturbance-adjustment decomposition is unchanged: $d\nu = (-J)^{-1}F\, d\tau$. Asymmetry therefore does not introduce a new economic channel. It changes the weights placed on the same channels. Two questions remain: where the bilateral threshold moves, and whether the aggregate NEER result survives.

\subsection{General asymmetric baselines}

Nothing in Section 5 used symmetry. At any balanced baseline, the system \eqref{eq:dp}--\eqref{eq:dtb} holds with the observed shares $(b_{ij}, h_i)$ and incomes $y_i$ as coefficients. The paper's appendix solves it for $N = 3$ at a general baseline, with unequal shares and sizes and existing tariffs. The bilateral slope against the bystander then has denominator $\det(-J) > 0$ and numerator quadratic in $\rho$:
\begin{equation}
\operatorname{sign} \frac{d\log e_{AC}}{d\tau}
= \operatorname{sign}\, \mathcal{N}(\rho),
\qquad
\mathcal{N}(\rho) = c_2\, \rho^2 + c_1\, \rho + c_0,
\qquad
c_2 = b_{AB}\, b_{AC}\, b_{CA}\, b_{CB}\, (1-h_A)(1-h_C)\, y_A\, y_C .
\label{eq:asym-quad}
\end{equation}
Two consequences:
\begin{itemize}
\item $c_2 > 0$ whenever $A$ and $C$ each buy from both foreign origins, since every share in the product is then positive. In that case $\mathcal{N}(\rho) > 0$ for $\rho$ large enough, so sufficiently strong cross-origin substitution reverses $e_{AC}$ at \emph{any} such baseline. Asymmetry moves the threshold; it cannot remove the reversal region.
\item The threshold itself is the largest root of $\mathcal{N}(\rho) = 0$. At the symmetric free-trade baseline that root is $3\rho_1^*$, recovering \eqref{eq:3c-slopes}.
\end{itemize}
The remaining coefficients $c_1, c_0$ depend on the full pattern of bilateral shares and sizes. Their interpretation is a reweighting of the symmetric channels. A country well positioned to replace $B$ in $A$'s import basket receives more of the diversion effect. A country heavily exposed to the fall in $B$'s income receives more of the offset. The sign of a bilateral response therefore depends on the trade network as well as on $(\eta, \rho)$, and no single universal threshold exists away from symmetry.

\subsection{Scaling exposure with $\kappa$}

To compare asymmetry across $N$, it is not useful to hold $B$'s absolute import share fixed, because the symmetric bilateral share itself shrinks like $1/(N-1)$. Instead, let the target $B$ carry a fixed multiple $\kappa$ of the symmetric bilateral share, with the weight matrix kept symmetric (hence balanced):
\begin{equation}
b_{AB} = b_{BA} = \frac{\kappa}{N-1},
\qquad
b_{AC} = b_{CA} = b_{BC} = b_{CB} = \frac{1 - \frac{\kappa}{N-1}}{N-2},
\qquad
b_{CC'} = \frac{1 - 2\,b_{CA}}{N-3} .
\label{eq:kappa-b}
\end{equation}
Here $\kappa = 1$ is exact symmetry, $\kappa > 1$ makes $B$ disproportionately important, and $\kappa < 1$ makes $B$ disproportionately small, in a way that remains comparable as $N$ changes. The bystanders are still interchangeable, so the reduced system again has two unknowns $(\nu_B, \nu_C)$ and solves in closed form (symbolically; replication package). With NEER weights $w_{Aj} = b_{Aj}$, the response $d\nu_A^E/d\tau$ is an explicit rational function of $(\rho, N, \kappa)$. It collapses to $-\rho_1^*/D_N$ exactly at $\kappa = 1$. Its numerator is $-\kappa\,\rho_1^*$ times a quadratic $Q(\rho)$ whose leading coefficient is proportional to $(\kappa - 1)(N - 1 - \kappa)$.

\subsection{Large-$N$ behavior}

\begin{result}[Asymmetric NEER]
Holding $\kappa$ fixed as $N \to \infty$,
\begin{equation}
N\,\frac{d\nu_A^E}{d\tau}
\;\longrightarrow\;
-\,\kappa\;\frac{\rho_1^*}{\rho + \alpha_D(\eta-1)}
\;=\; -\,\kappa \lim_{N\to\infty}\frac{N\rho_1^*}{D_N} :
\label{eq:kappa-neer}
\end{equation}
to leading order, asymmetric exposure scales the symmetric NEER response by the target's relative importance $\kappa$ and preserves its sign.
\end{result}

To leading order, asymmetry simply scales the symmetric effective appreciation by the relative importance of the tariff target. The symmetric result is therefore not economically knife-edge, even though its exact finite-$N$ cancellation is.

\subsection{When can the NEER sign reverse?}

At finite $N$ the sign of the exact solution follows from $Q$, and the direction of the asymmetry matters.
\begin{itemize}
\item $\kappa \geq 1$, a disproportionately \emph{large} target. All coefficients of $Q$ are positive for large $N$: the leading coefficient because $(\kappa-1)(N-1-\kappa) \geq 0$, and the other two because their leading terms are $N^2$ and $N^2\alpha_D(\eta-1)$. Then $d\nu_A^E/d\tau < 0$ for every $\rho$, and effective appreciation survives.
\item $\kappa < 1$, a disproportionately \emph{small} target. $Q$ has a single positive root,
\begin{equation}
\rho_E^*(\kappa, N) = \frac{N\rho_1^*}{1-\kappa}\,\big(1 + O(N^{-1})\big),
\label{eq:kappa-flip}
\end{equation}
and the NEER \emph{depreciates} for $\rho$ beyond it. The intuition follows directly from the aggregation argument. $A$ appreciates against $B$, but $B$ receives little weight in $A$'s NEER; if $A$ simultaneously depreciates against many bystanders, those movements dominate the average. As $\kappa \to 0$, the flip point approaches the bilateral threshold $N\rho_1^*$, because a weightless target leaves the NEER equal to the bystander rates.
\end{itemize}
The exact cancellation of \eqref{eq:N-neer} is therefore one-sided. Tariffs on disproportionately large partners preserve the conventional aggregate sign for every $\rho$. Tariffs on minor partners can reverse it, but only at $\rho$ of order $N$. The same dilution that makes individual bilateral reversals harder as $N$ grows also makes aggregate reversal under asymmetry difficult: such values lie far above the war threshold $\rho_1^*$ and, at moderate $N$, above estimated elasticities ($\rho_E^* \approx 22$ at $N = 10$, $\kappa = 0.5$, benchmark parameters). The symmetric result should be read as exact under symmetry and approximately robust more broadly, rather than as a knife-edge curiosity.

\section{Why the Canonical PTA Model Is Sharper}

The multilateral model predicts that an untreated exporter benefits from a discriminatory tariff only when cross-origin substitution is sufficiently strong. This may initially appear to conflict with the canonical theory of preferential trade agreements, in which discrimination against one exporter unambiguously improves the terms of trade of its competitor. The difference is not a competing mechanism. It is a difference in the set of margins permitted by the two environments.

\subsection{Competing exporters}

Countries $A$ (importer), $B$ and $C$ (exporters). Two goods: $x$ and numeraire $y$. Endowments: $A$ holds only $y$; $B$ and $C$ hold $x$ (and possibly $y$). Quasilinear preferences in each country,
\begin{equation}
U_i = u_i(x_i) + y_i, \qquad u_i' > 0,\; u_i'' < 0 .
\end{equation}
$A$ levies specific tariffs $t_B, t_C$ on imports of $x$ from $B$ and $C$. Let $q_j$ be the price received by exporter $j$ (the world price of its good) and $p$ the domestic price of $x$ in $A$. Both origins sell in $A$ at the same domestic price:
\begin{equation}
p = q_B + t_B = q_C + t_C .
\label{eq:pta-arb}
\end{equation}
From $u_i'(x_i) = (\text{local price of } x)$:
\begin{equation}
M_A(p) = x_A^d(p), \qquad
E_j(q_j) = \omega_j - x_j^d(q_j), \quad j = B, C,
\end{equation}
with $M_A' < 0$ and $E_j' > 0$. Market clearing for $x$:
\begin{equation}
M_A(p) = E_B(p - t_B) + E_C(p - t_C) .
\label{eq:pta-mc}
\end{equation}
Equation \eqref{eq:pta-mc} determines $p$ given $(t_B, t_C)$; then $q_B, q_C$ follow from \eqref{eq:pta-arb}. Totally differentiate \eqref{eq:pta-mc} with respect to $t_B$, holding $t_C$:
\begin{equation}
M_A'\, \frac{dp}{dt_B} = E_B'\Big( \frac{dp}{dt_B} - 1 \Big) + E_C'\, \frac{dp}{dt_B}
\quad\Longrightarrow\quad
\frac{dp}{dt_B} = \frac{E_B'}{E_B' + E_C' - M_A'} \in (0, 1) .
\end{equation}
Therefore
\begin{equation}
\frac{dq_B}{dt_B} = -\,\frac{E_C' - M_A'}{E_B' + E_C' - M_A'} \;<\; 0,
\qquad
\frac{dq_C}{dt_B} = \frac{dp}{dt_B} \;>\; 0 .
\end{equation}

\begin{result}[Viner--Mundell]
A discriminatory tariff on $B$ worsens $B$'s terms of trade and \emph{improves} $C$'s terms of trade. Equivalently, a preferential tariff cut toward $B$ (lower $t_B$) improves $B$'s and worsens $C$'s terms of trade.
\end{result}

\subsection{Which margins are missing?}

An improvement in $C$'s terms of trade is a rise in $e_{AC}$ \eqref{eq:tot}. The canonical model signs it unconditionally from $E_j' > 0$ and $M_A' < 0$; the nested model of Section 4 requires $\rho > \rho_3^* = 3\rho_1^*$. The canonical environment isolates foreign-supplier reallocation; the richer model embeds the same mechanism but allows several forces to oppose it. Read the threshold as the balance of three forces (benchmark values $\alpha_D = 0.8$, $\alpha_T = 0.4$, $\eta = 1.5$):
\begin{equation}
\rho_3^* = 3\rho_1^* = \underbrace{3}_{\text{$B$--$C$ linkage}}
\;+\; \underbrace{3\alpha_D(\eta-1)}_{\text{home switching}}
\;-\; \underbrace{3\alpha_T(1-\alpha_D)}_{\text{openness credit}}
\;=\; 3.0 + 1.2 - 0.24 = 3.96 .
\label{eq:pta-decomp}
\end{equation}

First, $B$ and $C$ trade with each other in the nested model. When $A$ tariffs $B$, the income lost in $B$ cuts $B$'s purchases from $C$, and the change in relative prices induces $C$ to substitute toward $B$'s cheapened goods. Both effects work against the simple idea that expenditure displaced from $B$ must move toward $C$; they are the source of the first term.

Second, consumers in $A$ can switch toward the domestic tradable rather than toward $C$. This home-switching channel strengthens with $\eta$ and with the home share $\alpha_D$, raising the $\rho$ required for diversion toward $C$ to dominate: the second term.

Third, openness works the other way. The import share $m = \alpha_T(1-\alpha_D)$ measures how much foreign expenditure the tariff can reallocate; a larger share lowers the threshold. This is the credit term.

The decomposition can be checked by recomputing the threshold with each margin removed (severing trade changes the baseline, so the pieces are approximately additive):

\begin{center}
\begin{tabular}{lll}
\toprule
restriction & threshold $\rho_3^*$ & piece removed \\
\midrule
none (benchmark) & $3.96$ & --- \\
no $B$--$C$ trade ($\beta_{BC} = \beta_{CB} = 0$) & $1.05$ & the linkage claim \\
no home tradables ($\alpha_D = 0$) & $1.80$ & the home-switching claim \\
no non-tradables ($\alpha_T = 1$) & $3.60$ & credit maximized (small alone) \\
$\alpha_D = 0$ and $\alpha_T = 1$ & none: reversal for all $\rho > 0$ & \\
no $B$--$C$ trade and $\alpha_D = 0$ & none: reversal for all $\rho > 0$ & \\
\bottomrule
\end{tabular}
\end{center}

Three readings:
\begin{itemize}
\item The largest piece is the $B$--$C$ linkage: severing it alone moves the threshold from $3.96$ to $1.05$.
\item Non-tradables are not an independent force: with the linkage severed and $\alpha_D = 0$, no threshold exists at any $\alpha_T$. The non-tradable share enters only through the credit.
\item Two routes reach $\rho_3^* \leq 0$: $\{\alpha_D = 0,\ \alpha_T = 1\}$ or $\{\alpha_D = 0,\ \beta_{BC} = \beta_{CB} = 0\}$. The canonical model takes both, adds quasilinearity, and in the Viner case sets $\rho = \infty$.
\end{itemize}

The canonical model removes most of these forces at once. $B$ and $C$ are competing exporters rather than full trading partners, so there is no $B$--$C$ demand linkage. $A$ does not produce the imported good, so displaced expenditure cannot move to a domestic variety. Quasilinearity removes the income effects. What remains is reallocation between the two foreign suppliers, signed directly by downward-sloping import demand and upward-sloping export supply.

\begin{result}
The canonical PTA model is the restriction of the three-country model to a subspace of structures on which $\rho_3^* \leq 0$. The unconditional Viner--Mundell sign is the threshold condition of Section 4 evaluated where it cannot bind. The nested model does not overturn the Viner--Mundell mechanism. It embeds that mechanism among the margins above, and $\rho > \rho_3^*$ states when substitution toward the untreated exporter dominates them.
\end{result}

\subsection{Relation to analytical and quantitative trade models}

Organized by which margins each model keeps or removes, the analytical literature maps onto the table:
\begin{itemize}
\item competing exporters (Bagwell--Staiger 1999; Ornelas 2005): $\beta_{BC} = \beta_{CB} = 0$, quasilinear;
\item symmetric endowment blocs (Krugman 1991; Bond--Syropoulos 1996): $\alpha_D \approx 0$, $\alpha_T = 1$, one elasticity;
\item Viner: $\rho = \infty$;
\item Kemp--Wan 1976: none of these restrictions, and no sign claims.
\end{itemize}

The applied literature sits on the other side of the same divide. Computable general equilibrium models of trade agreements in the GTAP tradition, and modern quantitative trade models, contain every force in \eqref{eq:pta-decomp}. They use a two-level import nest like \eqref{eq:mid}--\eqref{eq:inner}, full trade matrices, non-tradables, and income effects. In their simulations, a discriminatory tariff sometimes improves and sometimes worsens the excluded country's terms of trade, depending on the configuration. That is what crossing a threshold looks like when the threshold has not been written down. The models produce both signs, but no statement of when each occurs.

The division of labor is then simple. The analytical PTA models prove the sign by deleting the opposing forces. The applied models keep all the forces but report numbers, not conditions. The condition itself is $\rho > 3\rho_1^*$, and the contribution of the present framework is to make it explicit.

\section{Conclusion: Bilateral Reallocation Versus Aggregate Adjustment}

The familiar statement that an import tariff appreciates the tariffing country's currency is exact in a two-country world because there is only one foreign currency against which appreciation can occur. Once trade is multilateral, that statement combines two distinct economic adjustments. The first is aggregate substitution between home and foreign goods. The second is reallocation across foreign suppliers. A discriminatory tariff activates both. The aggregate margin pushes toward appreciation of the tariffing country, while the foreign-supplier margin redistributes that adjustment across trading partners and can cause individual bilateral exchange rates to move in the opposite direction.

The three-country model makes this distinction visible. $A$ always appreciates against the tariff target, but can depreciate against an untreated country when cross-origin substitution is sufficiently strong. Generalizing to $N$ countries strengthens the bilateral threshold because trade diversion is diluted across additional bystanders.

Aggregation produces a different result. Under symmetry, foreign-supplier reallocation cancels from the numerator of the nominal effective exchange-rate response. A unilateral tariff therefore appreciates the tariffing country's NEER for every $N$ and every cross-origin elasticity consistent with the multilateral Marshall--Lerner condition. The conventional two-country result thus survives multilateralization, but as an effective rather than universally bilateral statement.

The distinction also organizes the extensions. Broad protection treats foreign suppliers symmetrically and shuts down the reallocation margin. Reciprocal retaliation cancels the direct bilateral adjustment between the belligerents and can instead make foreign reallocation decisive for their effective exchange rates. Real exchange rates strengthen the aggregate appreciation through the tariff's direct effect on the home price level. Asymmetry breaks the exact aggregation result but largely preserves its economics. And the canonical Viner--Mundell model emerges as an environment in which the forces opposing foreign-supplier diversion have been removed; the threshold $\rho > 3\rho_1^*$ states when the diversion mechanism dominates once those forces are restored.

The common principle is simple:
\[
\boxed{\text{tariffs change both how much a country buys from abroad and from whom it buys it.}}
\]
The first margin governs the aggregate value of the currency. The second governs the distribution of that adjustment across bilateral exchange rates. Distinguishing the two clarifies both what survives from the two-country model and what genuinely changes in a multilateral world.

\end{document}
